% 本文件由 LaTeX 排版 Agent 根据有效 translation.json 自动生成。
% author=Athanasius; work_id=2812; title=致埃及与利比亚的主教们

\chapter{致\Place{埃及}与\Place{利比亚}的主教们}

% structure: p0001 source=h1 target=omitted reason=page-title-already-rendered-by-parent

% structure: p0002 source=h2 target=section reason=relative-heading-rank; base=h2

\section{第一章}

1. \Emph{基督警告祂的追随者提防假先知。}\par

我们的主和救主耶稣基督，正如\Person{路加}所写的，\BibleQuote{既有行动，又有训诲}\BibleRef{宗 1:1}，是在祂为我们的救恩显现之后成就的；因为祂来，正如\Person{若望}所说的，\BibleQuote{不是为审判世界，而是为使世界借着祂而得救}\BibleRef{若 3:17}。在其余事中，我们特别要赞叹祂的良善的这一实例：祂对于将要攻击我们的人并不缄默，而是预先明确告诉我们，好叫这些事发生时，我们能立时发现心意因祂的教导而坚定。因为祂说：\BibleQuote{将有假先知和假基督兴起，显大奇迹、大异事，以致连选民，如果可能，也要被欺骗。看，我预先告诉了你们}\BibleRef{玛 24:24}。祂在我们内贮蓄的恩宠的训诲和恩赐实在是多重的，超乎人的想象：如天上生活的模范，制伏魔鬼的能力，\OriginalTerm{义子}{adoption of sons}的名分，以及那极大且独特的恩宠，即认识圣父、认识圣言自身，以及圣神的恩赐。但人心极其倾向邪恶；此外，我们的仇敌魔鬼，嫉妒我们拥有如此大的福分，四处奔走，企图夺去那播在我们心中的圣言的种子。因此，好像要借着祂预言的警告，将祂的训诲视为祂自己特有的宝藏，印在我们心上，主说：\BibleQuote{你们要谨慎，不要受欺骗！因为将有许多人，假冒我的名字来说：我就是；又说：时期近了。你们切不可跟随他们}\BibleRef{路 21:8}。这是圣言赐予我们的厚礼，使我们不致被表象欺骗，反而无论这些事如何隐藏，我们都能借着圣神的恩宠越发辨别它们。因为邪恶的发明者兼大恶神，即魔鬼，本是极其可憎的，他一显露自己，就立刻被众人厌弃——如蛇，如龙，如咆哮的狮子，寻找可吞噬的人——所以他隐藏并掩饰自己的本相，狡猾地冒充那众人所渴慕的名号，好借着虚假的外表施行欺骗，从此将他所误导的人牢牢捆锁在自己的锁链里。又好比有人想趁父母不在时拐走他人的子女，便冒充他们父母的外貌，借此骗取孩子的亲情，将他们带到远方加以杀害；同样，这邪恶狡猾的魔鬼，对自己全无信心，又知道人们热爱真理，便冒充真理的外表，从而将自己的毒素散布于跟随他的人中。\par

2. \Emph{\Person{撒殚}伪装圣善，被基督徒识破。}\par

这样他欺骗了\Person{厄娃}，不说自己的话，却狡猾地采用天主的话，并曲解其意义。这样他向\Person{约伯}的妻子献恶计，劝她佯爱丈夫，同时教她亵渎天主。这狡猾的恶神正是如此用虚假的表演嘲弄人，迷惑各人并将其拖入自己邪恶的深坑。古时他欺骗了第一人\Person{亚当}，以为能借他使全人类臣服于己，便大为放肆，欢跃着说：\Quote{我的手寻到列国的财富，如同寻到鸟巢；如同人拾取遗下的蛋，我拾取了全地；无人能逃脱我，也无人能反抗我。}但主降临地上，敌人试探了祂的\OriginalTerm{救世工程}{human Economy}，却无法欺骗祂所取的肉躯，从那以后，那自许将占有全世界的，竟因祂的缘故连孩童都讥笑他：那骄傲者像麻雀一样被讥笑。如今婴孩将手放在毒蛇的穴口，嘲笑那欺骗了厄娃者；一切正确相信主的人，都将那说过\BibleQuote{我要升到云巅之上，与至高者同等}\BibleRef{依 14:14}的踏在脚下。这样他便受苦受辱；虽然他仍无耻地大胆伪装自己，但如今这可怜的家伙，反而更被那些额上有\OriginalTerm{印记}{Sign}的人识破；他们不但厌弃他，使他卑微，还使他蒙羞。因为他即使以爬蛇之形，化作\OriginalTerm{光明的天使}{angel of light}，他的欺骗也必无益；因为我们已受教导：\BibleQuote{即使是从天上来的天使，若传给你们与我们领受的不同的福音，他就当受\OriginalTerm{绝罚}{anathema}}\BibleRef{迦 1:8}。\par

3. 再者，他虽然隐藏其本性中的虚伪，并藉口唇假装说真话，然而\BibleQuote{我们并非不知道他的诡计}\BibleRef{格后 2:11}，反而能够用圣神论及他的话回答他：\BibleQuote{但天主对不信的人说：\InnerQuote{你为什么传述我的法令？}}又说：\BibleQuote{赞颂在罪人口中，是不相称的。}因为即使他说真话，这骗子也不配得到信任。圣经在叙述他在乐园中对厄娃的邪恶诡计时，已经指明了这一点；同样，主也谴责了他——首先在山中，当主揭露了\Quote{他胸铠的鳞片}，并指明这狡诈的神是谁，证实那试探他的不是某一位圣人，而是撒殚。因为主说：\BibleQuote{撒殚，退到我后面去！因为经上记载：\InnerQuote{你要朝拜上主你的天主，只应事奉祂。}}\BibleRef{玛 4:10}。再者，当主制止那些从坟墓向祂喊叫的魔鬼时，也是如此。因为尽管他们说的是真的，那时也没有说谎，说：\BibleQuote{你是天主子}和\BibleQuote{天主的圣者}\BibleRef{玛 8:29}；然而主不愿真理出自不洁之口，尤其是出自它们，以免他们以此为由，将其自身的毒计掺杂其中，并在人睡着时播种这些。因此，主不许他们说话，也不愿我们允许这种事，而是亲口吩咐我们说：\BibleQuote{你们要提防假先知！他们披着羊皮到你们这里来，内里却是凶残的豺狼。}\BibleRef{玛 7:15}，又藉着祂的圣宗徒的口说：\BibleQuote{不要凡神就信。}\BibleRef{若一 4:1}。这便是我们敌人的作为方式；一切异端的发明都与此类似，每一种异端都以魔鬼为其诡计之父，这魔鬼从起初就是杀人的，也是撒谎者。然而，他们羞于承认那可憎的名，便僭取我们救主那\BibleQuote{超越一切的名}\BibleRef{斐 2:9}，并以圣经的言词装饰自己，确是说那些话，却窃取了它们的真义；这样，他们用某种巧计伪装其虚妄的发明，也成了那些被他们引入歧途者的杀人者。\par

4. \Emph{只接受部分圣经而拒绝其他部分，毫无益处。}\par

因为\Person{马西昂}和\Person{摩尼}既拒绝法律，又从何接受福音呢？因为\WorkTitle{新约}源于\WorkTitle{旧约}，并为\WorkTitle{旧约}作证；如果他们拒绝\WorkTitle{旧约}，又怎能接受源于\WorkTitle{旧约}的呢？所以，\Person{保禄}是\WorkTitle{福音}的宗徒，\BibleQuote{这福音是天主先前藉着先知们在\WorkTitle{圣经}中所预许的}\BibleRef{罗 1:2}；我们的主亲自说：\BibleQuote{你们查考\WorkTitle{圣经}，因为正是这些为我作证}\BibleRef{若 5:39}。他们若不先查考为他所写的\WorkTitle{圣经}，又怎能承认主呢？门徒们说他们找到了他，\BibleQuote{\Person{梅瑟}在法律上和先知书上所记载的那一位}\BibleRef{若 1:45}。\OriginalTerm{撒杜塞人}{Sadducees}若不接受\WorkTitle{先知书}，法律对他们又算什么？因为赐下法律的天主，在法律中亲自应许也要兴起先知，所以同一位主既是法律的，也是先知的主；否认其中一个，也必须否认另一个。再者，如果\OriginalTerm{犹太人}{Jews}不承认按\WorkTitle{旧约}所期待降临的主，\WorkTitle{旧约}为他们又算什么？因为他们如果相信了\Person{梅瑟}的著作，也必相信主的话；因为他说：\BibleQuote{他写的是关于我}\BibleRef{若 5:46}。再者，对于拒绝圣言及其\OriginalTerm{降生成人的临在}{incarnate Presence}的\Place{萨莫萨塔}人，\WorkTitle{圣经}又算什么？这临在是\WorkTitle{旧约}和\WorkTitle{新约}所预示和宣告的。对于\OriginalTerm{亚略派}{Arians}，\WorkTitle{圣经}又有什么用，他们为何提出\WorkTitle{圣经}？这些人说圣言是受造物，并像\OriginalTerm{外邦人}{Gentiles}一样\BibleQuote{事奉受造之物，而不是事奉造物主}\BibleRef{罗 1:25}。因此，这些\OriginalTerm{异端}{heresies}中的每一个，就其各自捏造的独特不敬而言，都与\WorkTitle{圣经}毫不相干。他们的倡导者也明白这一点，即\WorkTitle{圣经}与每一个异端的教义都极为对立，甚至完全相反对；但为了欺骗更单纯的人（就像\BibleBook{箴}{言}所记的那些人：\BibleQuote{愚昧人凡话都信}\BibleRef{箴 14:15}），他们假装像他们的\BibleQuote{撒谎者的父亲魔鬼}\BibleRef{若 8:44} 一样研究并引用\WorkTitle{圣经}的言辞，以便借其言辞显得自己持有正确的信仰，从而说服他们那些可怜的追随者相信违背\WorkTitle{圣经}的道理。的确，在所有这些异端中，魔鬼就这样伪装了自己，并向他们暗示充满诡诈的话。主论到这些人说：\BibleQuote{将有\OriginalTerm{假基督}{false Christs}和\OriginalTerm{假先知}{false prophets}兴起，致使许多人受欺骗}\BibleRef{玛 24:24}。因此，魔鬼已来，借着各人说话，说：\Quote{我是基督，真理与我同在；}他使他们所有人都成了像他一样的说谎者。奇怪的是，尽管所有异端在各自所编造的邪恶发明上彼此分歧，却唯独在撒谎的共同目的上联合一致。因为他们有同一位父亲，将谎言的种子播撒在他们所有人里面。因此，忠信的基督徒和福音的真实门徒，因拥有恩宠能辨别属灵之事，并将信德的房屋建在磐石上，就坚定稳固，不受他们的欺骗。但单纯的人，如我之前所说的，对知识没有深入了解，这样的人只考虑所说的话，而不明白其含义，便立即被他们的诡计所掳去。所以，我们应当祈求并获得\OriginalTerm{辨别神类}{discerning spirits}的恩赐，以便每个人都能按照\Person{若望}的训示，知道应当拒绝谁，接纳谁为朋友和同信者。关于这些事，如果有人愿意详细说明，本当可以写下许多；因为异端的不敬和乖僻显得五花八门，欺骗者的诡计也极其可怕。然而，既然\WorkTitle{圣经}对我们已足够充分，所以我嘱咐那些愿意更深入了解这些事的人去诵读圣言，我现在便急切地向你们陈明那最值得注意的事，这也是我写这些的主要原因。\par

5. \Emph{亚略派试图以一个信经取代\OriginalTerm{尼西亚信经}{Nicene Creed}。}\par

我在这些地方逗留期间听闻（告知我的是真正正统的弟兄们），有些持亚略派观点的人聚集在一起，按自己的喜好起草了一份信仰宣认，并打算传话给你们，要你们要么签署他们所中意的，或者更准确地说，签署魔鬼所启发给他们的内容，要么拒绝签署就要遭受放逐。事实上，他们已经开始骚扰这些地区的主教们，由此清楚地暴露了他们的意图。因为他们制定这份文件，目的仅仅是为了施加放逐或其他惩罚，这种行为除了证明他们是基督徒的仇敌、魔鬼及其天使的同伙之外，还能是什么呢？尤其是他们违背那位仁慈的君主、我们最虔诚的\Person{君士坦提乌斯}皇帝的心意，四处散播他们所喜欢的。他们这样做非常狡猾，在我看来，主要怀有两个目的：第一，通过获得你们的签署，他们似乎可以消除\Person{亚略}这个名字所蒙受的恶名，并且他们自己可以躲过监视，仿佛并不信奉他的观点；再者，通过发布这些声明，他们可以给\Place{尼西亚}大公会议以及当时为反对亚略异端所颁布的信仰宣认蒙上阴影。然而，这种做法只是更明显地证明了他们自己的恶毒和异端思想。因为如果他们信仰正确，他们就会满足于整个普世公会议在\Place{尼西亚}所颁布的信仰宣认；如果他们觉得自己受到诽谤、被错误地称为亚略派，他们就不应该如此急切地去革新那些反对\Person{亚略}所写的内容，以免那些反对他的话也显得像是针对他们。然而，他们所走的并非这条路，他们是为自己而战，好像他们就是\Person{亚略}本人一样。看，他们完全无视真理，他们所说所做的一切都是为了亚略异端。他们竟敢质疑那些健全的信仰定义，擅自主张提出其他与之相反的教义，这不就是控告教父们，起来为那曾遭到他们反对和抗议的异端辩护吗？他们现在所写的，并非出于对真理的尊重，正如我先前说过的，而更像是出于嘲讽和诡计，为的是欺骗人；通过四处散发书信，他们可以让人们倾听这些想法，从而拖延他们被审判的时间；并通过将他们的不虔隐藏起来，使他们有空间扩展他们的异端，这异端\BibleQuote{就像毒疮一样}（见：\BibleRef{弟后 2:17}），四处侵蚀。\par

因此，他们扰乱并颠倒一切，然而即便如此，他们也不对自己的所作所为感到满足。因为他们每年都像要起草一份合同似的，聚在一起，假装撰写有关信仰的文件，这样一来，他们就更加让自己沦为笑柄和蒙受耻辱，因为他们的阐述不是被他人，而是被他们自己所拒绝。如果他们对自己先前的声明有任何信心，他们就不会想要起草别的；再者，离开了这最后的声明，他们现在也不会记下目前这份，而毫无疑问，按照他们的惯常做法，他们很快又会修改它，只要他们找到借口，对他们惯常密谋反对的某些人采取行动。因为当他们设计陷害某人时，他们就会大张旗鼓地撰写关于信仰的文件；正如\Person{比拉多}洗手一样，他们便想通过撰写文件来毁灭那些正确信仰基督的人，希望藉着对信仰作出定义，如我多次说过的，显得自己与错谬教义的指控无关。但他们无法隐藏自己，也无法逃脱；因为甚至在他们为自己辩护时，他们也在不断控告自己。这是理所当然的，因为对于拿出证据反对他们的人，他们不但不去回应，反而只是说服自己去相信自己想要的任何东西。当罪犯自己成为自己的法官时，怎能获得开释呢？因此，他们不断地撰写，又不断地更改自己先前的声明，从而显露出一种动摇的信德，或者更确切地说，一种昭然若揭的不信与乖僻。依我看来，他们必然落入这般下场；因为既然他们背离了真理，又想推翻在\Place{尼西亚}所制定的那健全的信仰宣认，用圣经上的话来说，他们\BibleQuote{喜爱流浪，没有约束自己的脚}（见：\BibleRef{耶 14:10}）；因此，就像古时的\Place{耶路撒冷}一样，他们在变化中劳苦经营，时而写这个，时而写那个，无非是为了拖延时间，好让他们继续与基督为敌，继续做人类的欺骗者。\par

\Emph{阿卡基乌斯派实际上是亚略派}\par

那么，谁若真的关心真理，还会愿意容忍这些人呢？谁不会理所应当地拒绝他们的作品？谁不会谴责他们的大胆妄为：他们人数极少，却要让自己的决定凌驾于一切之上，并且贪求自己在角落里举行、情形可疑的会议的至高地位，要强行取消那纯洁无玷的\OriginalTerm{大公会议}{Ecumenic Council}的决议？这些人因拥护这\OriginalTerm{敌基督的异端}{Antichristian heresy}而受到\Person{优西比乌}及其同党的提拔，竟敢擅自界定信理，而他们本该如\Person{盖法}那样，像罪犯一样受审，却反倒自任判官。他们编写了一部\OriginalTerm{\WorkTitle{塔利亚}}{Thalia}，竟想让人奉之为信德的标准，而他们自己甚至还没有确定信什么。谁不知道，\Person{培恩塔颇里的塞康杜斯}，虽早已多次被贬黜，却因\OriginalTerm{亚略的疯狂}{Arian madness}而被他们接纳；而如今在\Place{劳狄刻雅}的\Person{乔治}、太监\Person{莱翁提乌斯}，以及在他之前的\Person{斯特凡努斯}和\Place{赫拉克勒亚}的\Person{德奥多罗}，都是由他们提拔的？\Person{乌尔撒奇乌斯}和\Person{瓦伦斯}也一样，他们年轻时最初受教于\Person{亚略}，虽从前已从\OriginalTerm{司铎品}{Priesthood}中被贬黜，后来却因不虔敬而获得主教头衔；\Person{阿卡奇乌斯}、\Person{帕特罗斐鲁斯}和\Person{纳尔奇索斯}也是如此，他们曾在各种不虔敬的事上最为积极。这些人在\Place{撒底卡}的\OriginalTerm{大公会议}{great Synod}上被贬黜；而如今在\Place{色巴斯特}的\Person{优斯塔提乌斯}、\Person{德莫斐鲁斯}和\Person{格尔米纽斯}、\Person{优多克西乌斯}和\Person{巴西略}，这些那不虔敬之事的支持者，也是以同样方式被擢升的。至于\Person{凯克罗皮乌斯}、那个他们称为\Person{奥克森提乌斯}的人，以及骗子\Person{埃皮克泰图斯}，我无需多言，因为众人皆知他们是以何种方式、藉何种借口、被哪些我们的敌人所提拔，为的是让他们能对自己阴谋加害的正统主教进行诬告。因为，尽管他们住在\OriginalTerm{八十邮站}{eighty posts}之外，不为百姓所知，却凭着自己的不虔敬，为自己买到了主教头衔。出于同样原因，他们如今又雇佣了一个\Place{卡帕多细亚}的\Person{乔治}，企图强加给你们。但对他，和其他人一样，无需表示任何尊重；因为这一带有传闻说，他甚至不是基督徒，而是沉迷于\OriginalTerm{偶像崇拜}{worship of idols}；而且他有着\OriginalTerm{刽子手般的性情}{hangman's temper}。他们已把这样一个人，正如所描述的那样，纳入自己的队伍，以便能够伤害、掠夺和杀戮；因为他在这些事上大有本事，却对基督徒信仰的\OriginalTerm{基本教义}{very principles of the Christian faith}一无所知。\par

8. \Emph{出自恶人的话语，即使是圣经上的，也是坏的。}\par

这些就是这些人反对真理的阴谋诡计：但他们的图谋已显露在全世界面前，尽管他们千方百计，如同\OriginalTerm{鳗鱼}{eels}，试图滑脱掌握，逃避被认出是\OriginalTerm{基督的仇敌}{enemies of Christ}。因此，我恳求你们，你们中间不要有人受骗，不要有人被他们诱惑；相反，当思量有一种\OriginalTerm{犹太式的不虔敬}{judaical impiety}正侵入基督徒信仰，你们众人都要为主大发热心；每一个人都要坚守我们从\OriginalTerm{教父}{Fathers}所领受的信德，那是在\Place{尼西亚}聚集的人们以书面记录下来的，并要拒绝那些试图在这信仰上标新立异的人。无论他们写出怎样出自圣经的词句，都不要容忍他们的作品；无论他们怎样说正统的言辞，也不要听信他们所说的；因为他们并非心口如一，而是披上这些言辞，犹如\OriginalTerm{羊皮}{sheeps' clothing}，内心却随从\Person{亚略}，效法那一切\OriginalTerm{异端}{heresies}的创始者\OriginalTerm{魔鬼}{the devil}的行径。因为这魔鬼也曾引用圣经的话，却被我们的\OriginalTerm{救主}{Saviour}驳得哑口无言；因为倘若他引用时真正是那样理解，他就不会从天堂坠落了；但如今，他既因骄傲而堕落，就在言语中狡猾地掩饰，并时常恶意地试图用\OriginalTerm{外邦人}{Gentiles}的诡辩和巧言，引诱人走入歧途。假如他们的这些阐释是出自正统人士，出自像伟大的\OriginalTerm{宣信者}{Confessor}\Person{何西乌斯}，\Place{高卢}的\Person{马克西米努斯}或其继任者，或像东方的主教\Person{菲洛戈尼乌斯}和\Person{优斯塔提乌斯}，\Place{罗马}的\Person{犹利乌斯}和\Person{利伯略}，\Place{默西亚}的\Person{齐里亚科}，\Place{希腊}的\Person{匹斯图斯}和\Person{阿里斯特乌斯}，\Place{达契亚}的\Person{西尔维斯特}和\Person{普罗托格内斯}，\Place{卡帕多细亚}的\Person{莱翁提乌斯}和\Person{优普塞基乌斯}，\Place{阿非利加}的\Person{凯奇利安}，\Place{意大利}的\Person{优斯托尔吉乌斯}，\Place{西西里}的\Person{卡皮托}，\Place{耶路撒冷}的\Person{玛加略}，\Place{君士坦丁堡}的\Person{亚历山大}，\Place{赫拉克勒亚}的\Person{派德罗斯}，或者像那些伟大的主教，即来自\Place{亚美尼亚}和\Place{本都}的\Person{梅勒提乌斯}、\Person{巴西略}和\Person{隆吉安}以及其余人，来自\Place{基利家}的\Person{卢普斯}和\Person{安菲翁}，来自\Place{美索不达米亚}的\Person{雅各伯}和其余人，还有我们可敬的\Person{亚历山大}，以及其他与他们持相同意见的人——那么他们的陈述中就绝不会有什么可疑之处，因为正人君子的品性\OriginalTerm{真诚无欺}{sincere and incapable of fraud}。\par

9. \Emph{因为这样的话不过是他们的伪装。}\par

他们出自那些受雇为异端辩护的人，并且正如天主的箴言所说：\BibleQuote{恶人的言语是埋伏等候}，\BibleQuote{恶人的口倾吐邪恶}，\BibleQuote{恶人的计谋是欺诈}；弟兄们，正如主所说的，我们应当警醒谨慎，免得有人用花言巧语和诡诈欺骗我们；免得有人前来假装说\Quote{我宣讲基督}，不久之后却发现他是\OriginalTerm{反基督}{Antichrist}。凡是为\OriginalTerm{亚略}{Arian}的疯狂而来你们这里的，这些人就是\OriginalTerm{反基督}{Antichrist}。因为在你们中间有什么缺失，竟需要有人从外面来你们这里呢？或者\Place{埃及}、\Place{利比亚}和\Place{亚历山大里亚}的各教会，有什么如此急需，竟让这些人用金钱购买\OriginalTerm{主教职}{Episcopate}，而非木材货物，并侵入不属于他们的教会呢？谁不知道，谁不清晰看出，他们做这一切都是为了支持他们的不虔敬呢？因此，即使他们使自己哑口无言，或是在衣边上佩戴比法利塞人更宽的衣穗，长篇大论，练习声调，也不应相信他们；因为推荐忠信基督徒的，不是讲话的方式，而是内心的意向和虔敬的生活。因此，撒杜塞人和黑落德党人，虽然他们口中常有法律，却被我们的救主斥责，救主对他们说：\BibleQuote{你们错了，不明了圣经，也不明了天主的德能}\BibleRef{玛 22:29}；众人都见证了那些假装引用法律言辞的人被揭露，他们在心思上是异端者和天主的仇敌。他们用这些言词确实欺骗了别人，但当我们的主降生成人时，他们却无法欺骗他；\BibleQuote{因为圣言成了血肉}，他\BibleQuote{认识人的思想无非都是虚幻}。他就这样揭露了犹太人的挑剔，说：\BibleQuote{如果天主是你们的父亲，你们必爱我，因为我是由圣父出发而来到你们这里的。}如今，这些人似乎同样行事；因为他们掩饰自己的真实想法，然后利用圣经的语言写文章，将其作为诱饵，勾引无知的人，把他们诱入自己的邪恶之中。\par

10. \Emph{如果他们想被听取，就应先谴责\OriginalTerm{亚略}{Arius}。}\par

请想想，是否确是如此。如果他们在没有理由的情况下，撰写信仰宣认，这便是多余，或许还是有害之举，因为在无人追问之际，他们却为言辞争论提供机会，扰乱弟兄们的纯朴之心，在他们中间散布从未想到过的观念。如果他们试图为自己就亚略异端撰写辩护，他们本应先铲除那些已经生发的邪恶种子，并谴责那些制造它们的人，然后以正确的声明取代先前的说法；否则就让他们公开为\OriginalTerm{亚略}{Arius}的意见辩护，这样他们就不再是暗中，而是公开地自显为基督的仇敌，使众人躲避他们如同躲避蛇的面目。但现在他们却隐藏那些意见，借口撰写别的事情；恰如一位外科医生，被召来照顾一个受伤且痛苦的人，进来后却对他的伤口只字不提，反而谈论他那健全的肢体。这种人无疑是最为愚蠢的，因为对他被请来的事情不提，反而去谈论那些无人需要的事情。同样，这些人略去关乎异端的事，却着手撰写其他主题；可是，如果他们对信仰有任何关心，或者对基督有任何爱，他们就应当首先清除那些反对他的亵渎言语，然后再以正确的言词说话和写作。但他们自己既不这样做，也不允许那些渴望这样做的人去做，这或是出于无知，或是出于诡诈和诡计。\par

11. \Emph{若在一方面行恶，在另一方面行善，便毫无利益。}\par

若他们出于无知而这样做，必被指为鲁莽，因为他们对自己所不知的事竟擅作断定；然而，倘若他们是明知而掩饰，则其罪更大，因为他们在图谋个人利益上无所不察，但在书写关于我主的信德时，却加以戏弄，宁可做任何事也不讲真话；他们隐瞒那些使他们的异端受指摘的细节，而只提出圣经上的话。这显然是窃取真理的行径，充满一切不义的勾当；故此，我确信你的敬虔之心从下述实例中不难看出这一点。一个人若被控犯奸淫，不会以未犯盗窃来为自己辩护清白；同样，若有人提出杀人的指控，也不会容许被告说：\Quote{我们并未发虚誓，却保存了所托付于我们的寄托物}来自辩。这只不过是儿戏，而不是对指控的辩驳和对真理的展示。因为杀人同寄托物，奸淫同盗窃，有何相干？这些恶行固然都出自同一内心，彼此关联；但在驳斥某一被控罪行时，却互不相干。因此，正如\BibleBook{若苏厄}{书}所载，当\Person{阿干}被控犯有偷窃时，他并未以自己热心战事为借口开脱，而是被定罪后遭全体民众用石头砸死。又如\Person{撒乌耳}被控疏忽和违犯律法时，他提出其他事例也未能使其诉讼获益。\EditorialNote{见：\BibleRef{撒上 15}} 因为关于一项指控的辩护，不能用来获得另一项指控的无罪开释；但万事若按法律和正义而行，一个人必须就他被指控的各项进行辩护，要么驳倒既成事实，要么承认下来，并许诺改正，不再重犯。然而，他若犯了罪却不肯承认，为掩盖真相，避而不谈所问之事，却另讲别事，便清楚表明他行事不端，而且自知有罪。其实，何必多言，既然这些人本身就是亚略异端的指控者？他们既不敢明说，反将其亵渎之词隐藏起来，这就显明他们深知此异端是与真理分隔、格格不入的。然而，既然他们自己将它隐藏，不敢明言，便必须由我来揭去他们不信不虔的面纱，将这异端公然陈列，因为我晓得\Person{亚略}及其同党当初的言论，也晓得他们是从教会中被驱逐，并从圣秩中罢免的。但在此我先请求宽恕，因我将引述那些污秽的话；我使用它们，非因我如此认为，而是为揭发这些异端者。\par

% structure: p0023 source=h2 target=section reason=relative-heading-rank; base=h2

\section{第二章}

12. \Emph{亚略派的言论}\par

蒙主恩召的\Person{亚历山大}主教，因\Person{亚略}持有并坚持下列见解而将他逐出教会：\Quote{天主并非永远是圣父；圣子并非永远存在；而是，既然万物都从无中造成，圣子也是从无中造成的；既知万物皆受造，他也是受造物，且是被造之物；既知万物曾一时不存在，后来又造成，曾有一时间，天主圣言本身也不存在；他在受生之前并不存在，却有存在之始：因为他是在天主决定产生他时才产生的；因为他也是天主的其余化工之一。又因他本性可变，他之所以继续为善，只因他凭自己的自由意志选择；所以，凡他愿意时，他便同其他万物一样，能够改变。因此，天主预知他将为善，预先赐予他那后来他因德行才能获得的光荣；他现在是因天主预先知道的那些行为而成为善的。 所以，他们说，基督并非真天主，而是因分沾天主性体而得称为天主，正如其他所有受造物一样。他们又说，他不是那本性在圣父内、且为圣父性体所专有的圣言，也不是圣父创造此世所使用的固有智慧；而是另有一位圣言，专属于圣父，另有一位智慧，专属于圣父，圣父用这智慧造成了此圣言。他们说，主自己在概念上，就具有理性之物而言，被称为\OriginalTerm{圣言（理性）}{Word (Reason)}；在概念上，就具有智慧之物而言，被称为\OriginalTerm{智慧}{Wisdom}。 实际上，他们说，如同万物在本质上与圣父分离且迥异，他也在各方面与圣父的本质分离且迥异，本质上属于受造和造成之物，且为其中之一；因为他是一个受造物、一件被造物、一项工程。此外，他们说，天主不是为了他自己而造了我们，却是为了我们的缘故而造了他。因为他们说：\InnerQuote{天主本是独一的，圣言并未与他同在，但后来，当他愿意生产我们时，他才造了圣言；自从他被造之时起，便称他为圣言、子及智慧，好借着他创造我们。正如万物由天主的意志而存在，先前并不存在；同样，他也是由天主的意志而造，先前并不存在。因为圣言并非父所固有、本然的苗裔，而是由恩宠产生的：因为自有的天主，借着自己的意志，造了那原本不存在的子，天主也是用此意志造了万物，使之产生、受造、并愿意它们存在。} 他们还说，基督并非天主本然和真正的能力；而是如同蝗虫和蚱蜢被称为一种\InnerQuote{能力}，他同样也因之被称为圣父的\InnerQuote{能力}。此外，他还说，父对子而言是隐秘的，子既不能完全、准确地看见父，也不能认识父。因为，既有存在之始，他便不能认识那无始者；他所认识、所看见的，是按照他自身的度量，依比例而认识、看见，正如我们也按自己的能力认识、看见一样。他又补充说，子不仅不完全认识自己的父，甚至连自己的性体也不认识。}\par

13. \Emph{由圣经反驳亚略言论的论证}\par

因为坚持这些及类似的主张，\Person{亚略}被宣布为异端者；就我而言，虽然我只是在记录它们，我通过思考相反的教理，并持守真信仰的意识，已洁净自己。因为所有从各地聚集在\Place{尼西亚}大公会议的主教们，一听到这些声明便掩耳；所有人同声谴责这异端，并将其绝罚，宣布它与教会的信仰相异且疏离。并非强迫引导审判官们作出这个决定，而是他们全都深思熟虑地维护了真理；他们做得公正且正确。因为不信通过这些人在传播，或更好说，是一种与圣经相悖的犹太教思想，其上附著外邦人的迷信，以致持有这些主张的人甚至不能再被称为基督徒，因为它们全都与圣经相悖。例如，\Person{若望}说：\BibleQuote{在起初已有圣言} \BibleRef{若 1:1}；但这些人却说：\Quote{祂在受生之前并不存在。}他又写道：\BibleQuote{我们也在那真实者内，即在他的圣子耶稣基督内；他就是真天主，和永生} \BibleRef{若一 5:20}；但这些人，如同在反驳此点，声称基督不是真天主，而只是像其他受造物一样，因分享天主性而被称作天主。宗徒责备外邦人，因为他们崇拜受造物，说：\BibleQuote{他们敬拜事奉受造物，代替造物主。}但如果这些人说主是受造物，并把他当作受造物来崇拜，他们与外邦人有何区别？如果他们持有这种主张，这节经文岂不也反对他们；有福的\Person{保禄}岂不正是在责备他们吗？主也说：\BibleQuote{我与我父是一体}；又说：\BibleQuote{谁看见我，就是看见了父}；而且他所派遣宣讲的宗徒写道：\BibleQuote他是天主光荣的辉耀，和他本体的真像 \BibleRef{希 1:3}。但这些人竟敢分离他们，并说祂与圣父的\OriginalTerm{本质}{essence}和永恒是疏离的；还邪恶地认为祂是可变的，意识不到这样说，他们就使祂不与圣父合一，而与受造物合一。谁不明白，辉耀不能与光分离，而是本性上属于光，与光共存，并非在光之后产生？再者，当圣父说：\BibleQuote{这是我的爱子} \BibleRef{玛 17:5}，当圣经说祂是圣父的\BibleQuote{圣言}，\BibleQuote{藉着祂诸天被建立}，简言之，\BibleQuote{万有都是藉着祂而造成的} \BibleRef{若 1:3}；这些新教义和虚构的发明者却声称，另有圣父的圣言，另有圣父的智慧，而祂只是因为受造物中有理性者才被概念地称作圣言和智慧，却未察觉到这有何等荒谬。\par

14. \Emph{从圣经反驳亚略派陈述的论据。}\par

但如果他被称作圣言和智慧是由于我们而有的虚构，那么他究竟是什么他们却不能说明。因为如果圣经断言主是这两者，而这些人却不允许他如此，显然他们因对圣经的无神反对而完全否认他的存在。信友们能从圣父他自己的声音，从崇拜他的天使，以及从记载关于他的圣人们得出结论；但是这些人，因没有纯洁的心灵，且不能忍受聆听教导天主的属神之人的话语，或许甚至能从与他们相似的魔鬼那里学到一些东西，因为魔鬼说到他，并非好像有许多别的存在，而是，如同认识他唯一一样，说：\BibleQuote{你是天主的圣者}，和\BibleQuote{天主子} \BibleRef{谷 1:24}。那位在山上试探他、向他们暗示此异端者，也没有说\Quote{如果你也是天主子}，仿佛除了他还有其他，而是说\Quote{如果你是天主子}，如同他是唯一的那位。但正如外邦人，从唯一神的观念堕落，陷入多神论，这些可惊的人也一样，不信圣父的圣言是唯一的，而接受了多个圣言的理念，并且否认那位真是天主和真圣言者，竟敢把他设想为受造物，未察觉这想法充满何等的邪恶。因为如果他是受造物，他又如何同时是受造物的造物主？又如何是子和智慧和圣言？因为圣言不是受造，而是受生；受造物不是子，而是产物。如果所有受造物都是藉着他而造成的，而他也是受造物，那么他是由谁而造成的呢？被造的必须透过某一位而产生；事实上正是透过圣言而产生；因为他自己不是受造物，而是圣父的圣言。再者，如果在圣父内除了主之外另有智慧，那么智慧就源于智慧了；如果天主的圣言是天主的智慧，那么圣言就源于圣言了；如果子是圣父的圣言，那么子就必须是在子内造成的了。\par

15. \Emph{从圣经反驳亚略派陈述的论据。}\par

主既然说了：\BibleQuote{我在父内，父也在我内}\BibleRef{若 14:10}，怎么能在圣父内还有另一位，而主自己也借着他而受造呢？若望怎么越过了那另一位，只论到这一位说：\BibleQuote{万物是藉着他而造成的；凡受造的，没有一样不是藉着他而造成的}\BibleRef{若 1:3}呢？如果万物是藉着他而造成的，是按天主的旨意造成的，他自己怎能是受造物之一呢？还有，当宗徒说：\BibleQuote{万物的终向和万物根源的}\BibleRef{希 2:10}，这些人怎么能说，我们不是为他而受造，反倒是他为我们而受造呢？如果真是这样，他就应该说：\Quote{圣言为谁而受造}；但他没有这样说，而是说：\BibleQuote{万物的终向和万物根源的}，这就证明这些人是异端者和说谎者。此外，他们既然胆敢说在天主内还有另一位\OriginalTerm{圣言}{Word}，又无法从圣经中提出任何明确证据，那就让他们至少指出他的一项工程，或圣父在没有这位圣言的情况下所作的一项工程；这样他们的想法至少能显得有点根据。真实圣言的工程对众人是明显的，从而可以藉着类比从这些工程中观看到他。因为正如我们见了受造界，就设想天主是它的创造者；同样，当我们看到其中没有一样是无秩序的，万物都在秩序和\OriginalTerm{天主眷顾}{providence}中运行和持续，我们就推断有一位天主的圣言，他统御万有，管理万有。圣经也这样作证，宣布他是天主的圣言，说：\BibleQuote{万物是藉着他而造成的；凡受造的，没有一样不是藉着他而造成的}\BibleRef{若 1:3}。至于他们所讲的那另一位圣言，他们却拿不出任何言语或工程作为证据。不，甚至圣父亲自说：\BibleQuote{这是我的爱子}\BibleRef{玛 17:5}时，也表明除他以外再没有别的。\par

16. \Emph{亚略派与摩尼派的相似之处}\par

因此，就这些教义而言，这些奇妙的人物现在已与摩尼派同流合污了。因为摩尼派也承认有一位善的天主存在，但这只停留在名称上，他们却指不出他任何可见或不可见的工程。但他们既然否定那位真实而确实是天主的、天地和一切不可见之物的创造者，他们就只是捏造神话的人。我看这些心术不正的人正是如此。他们看见了那独在圣父内的真实圣言的工程，却否定他，为自己另造了一位圣言，而这位圣言的存在，他们既不能从他的工程，也不能从他人的证言加以证明。除非他们对天主怀有一种神话般的观念，认为他像人一样是个复合的存在，说话后又改变言语，像人一样运用理智和智慧；他们不知这样的意见使他们陷入何等荒谬的地步。因为如果天主有一连串的话语，他们定然得把他当作一个人。如果那些话语从他发出，随后又消失，他们便犯了更大的不敬之罪，因为他们使由自有者天主所发出的归于虚无。如果他们设想天主有所生，那么肯定更好、更虔诚的说法是：他生了一位圣言，这位圣言是他\OriginalTerm{天主性}{Godhead}的圆满，在他内蕴藏着一切智慧的宝藏，他与圣父永远同在，万物都是藉着他而造成的；而不是设想天主是许多无处可寻的话语之父，或者把那本性单纯的天主描绘成由众多部分组合而成，且受制于人的情欲，变化无常。再者，当宗徒说：\BibleQuote{基督是天主的德能和天主的智慧}\BibleRef{格前 1:24}，这些人竟把他当作众多能力中的一种；不，比这更糟，他们身为罪犯，竟把他比作天主为惩罚人而派遣的蝗虫和其他无灵之物。再者，当主说：\BibleQuote{除了子和子所愿意启示的人外，也没有人认识父}\BibleRef{玛 11:27}；又说：\BibleQuote{这不是说有人看见过父，只有那从天主来的，才看见过父}\BibleRef{若 6:46}时，那些说圣子不能完全看见圣父、也不能完全认识圣父的人，岂不真成了天主的仇敌吗？如果主说：\BibleQuote{正如父认识我，我也认识父一样}\BibleRef{若 10:15}，而圣父既不是部分地认识圣子，那么这些人胡言乱语，说圣子只部分地认识圣父，而非完全认识，岂不是疯了吗？再者，如果圣子有存在的开端，而万物也都有开端，让他们说说，哪个在先。其实他们无话可说，即便用尽手段，也无法证明圣言的这种开端。因为他乃是圣父真正而本然的圣子，且\BibleQuote{在起初已有圣言，圣言与天主同在，圣言就是天主}\BibleRef{若 1:1}。至于他们断言圣子不认识自己的本体，除了谴责他们的狂妄以外，再予答复便是多余了；因为圣言既将关于圣父和他自己的知识赐给众人，又责备那些不认识自己的人，他怎能不认识自己呢？\par

17. \Emph{从圣经引证反对亚略派的说法}\par

但他们是：经上记着：\BibleQuote{主在造化之初就创造了我，作为他工程的开端。} 无知的愚昧人啊！经上也称他为\BibleQuote{仆人}，\BibleQuote{婢女的儿子}，\BibleQuote{羔羊}，\BibleQuote{羊}，并且记载他劳苦、干渴、受鞭打、受苦痛。但圣经中这样描述他，自有其合理的根据和原因，这是因为他成了人，成了人子，取了奴仆的形体，就是人的血肉：因为\Person{若望}说：\BibleQuote{圣言成了血肉}（\BibleRef{若 1:14}）。既然他成了人，就不该因为这些表述而反感；因为受造、诞生、成形、受劳苦、受痛苦、死亡并从死者中复活，都是人性所固有的。正如他是圣父的\OriginalTerm{圣言}{Word}和\OriginalTerm{智慧}{Wisdom}，具有圣父的一切属性，他的永恒，他的不变性，在一切方面和一切事上与圣父相似，既无先后，又与圣父同性同体，他正是\OriginalTerm{天主性}{Godhead}的真像，是\OriginalTerm{造物主}{Creator}，而非\OriginalTerm{受造物}{creature}——（因为他在本质上与圣父相同，就不可能是受造物，而是造物主，正如他自己所说的：\BibleQuote{我父到现在一直工作，我也工作}（\BibleRef{若 5:17}）——同样，当他成了人，取了我们血肉时，就必定有\InnerQuote{受造}、\InnerQuote{被造}的说法，这是所有血肉之躯的特质；然而这些人，就像犹太酒贩子在酒里掺水一样，贬低圣言，将他的天主性屈从于他们关于受造物的观念。因此，\OriginalTerm{教父}{Fathers}们理所当然且公义地愤怒了，绝罚了这极不虔敬的\OriginalTerm{异端}{heresy}；如今这些人对此小心翼翼，不敢声张，因为这种异端容易驳斥，且处处站不住脚。我这里所陈述的，只不过是用来驳斥他们学说的少数论证；但如果有人希望更深入地反驳他们，就会发现这种异端与\OriginalTerm{异教}{heathenism}相去无几，而且是所有异端中最卑下、最渣滓的。其他异端或在主的身体或降生成人方面犯错误，扭曲真理，各有各的方式；或者完全否认主曾寓居人世，就像犹太人错误地想象的那样。唯独这种异端比其余的更狂傲，竟敢攻击天主性本身，断言圣言根本不存在，断言圣父并非永远是父；因此，人们可以合理地说，下面这首圣咏就是针对他们写的：\BibleQuote{愚人心中说：没有天主。他们都已败坏，行为可憎。}\par

18. \Emph{如果\OriginalTerm{亚略派}{Arians}认为自己是正确的，他们就会公开明说。}\par

\Quote{但是，}他们说，\Quote{我们强壮有力，能用许多诡计捍卫我们的异端。}他们若真有纯朴的\OriginalTerm{信德}{faith}，而非诡计和外邦人的诡辩术来捍卫，那倒还算有个较有理的回答。然而，如果他们有信心，并知道这异端符合教会教义，就让他们公开表明观点吧；因为没有人点了灯，放在斗底下（\BibleRef{玛 5:15}），而是放在灯台上，照耀所有进来的人。因此，如果他们真能捍卫，就让他们把上面指摘给他们的那些观点写成文字，像摆出灯一样，将他们的异端赤裸裸地展示在众人眼前，并公开指控真福忆念的\Person{亚历山大}主教，说他因\Person{亚略}宣认这些观点而不公正地逐出了他；还要责怪\Place{尼西亚}大公会议，因为它颁布了真实信仰的书面宣信，代替他们的不虔。但我敢肯定，他们不会这样做，因为他们对自己所捏造并图谋散布的邪说之恶并非毫无所知；他们非常明白，即使起初能用虚妄的欺诈迷惑单纯的人，这些妄想也必很快熄灭，\BibleQuote{如同恶人的光}（\BibleRef{约 18:5}）亦必熄灭，而他们自己也必到处被贴上真理仇敌的烙印。因此，尽管他们行事尽皆愚昧，说话如同傻子，但至少在这件事上，他们倒像\BibleQuote{今世之子}（\BibleRef{路 16:8}）那样聪明，把灯藏在斗底下，好让人以为它在发光，免得一旦显露就被定罪熄灭。就这样，当这异端的创始人\Person{亚略}，\Person{优西比乌}的同党，经由\Person{优西比乌}及其党羽的促请，被召去见真福忆念的\Person{君士坦丁}·奥古斯都，并被要求书面宣示自己的信仰时，这个狡猾的家伙写了一份，却隐去了他不虔信理的特有词句，模仿魔鬼的伎俩，只引圣经上的简朴言词，照本宣科。真福\Person{君士坦丁}对他说：\Quote{如果你心中没有别样思想，除了这些以外，就让真理为你作证；你若发虚誓，主必惩罚你。}这个不幸的人就发誓说他绝无别样思想，也从未说过或想过与自己书面所写相左的东西。可是，他一出去，就倒下了，仿佛为自己的罪行偿付了惩罚，他\BibleQuote{头朝下跌倒，腹部崩裂}（\BibleRef{宗 1:18}）。\par

19. \Emph{\Person{亚略}之死的意义。}\par

死亡固然是众人的共同结局，我们不应侮辱死者，即使他是仇敌，因为同一件事或许在傍晚之前就临到我们身上，这并不确定。但\Person{亚略}的结局并非寻常方式，因此值得叙述。\Person{尤西比乌}及其同伙威胁要将他带入教会，\Place{君士坦丁堡}主教\Person{亚历山大}抵抗他们；但\Person{亚略}倚仗\Person{尤西比乌}的暴力和恫吓。那天是安息日，他指望第二天能参加共融。因此他们之间发生了激烈的争斗；一方在威胁，\Person{亚历山大}在祈祷。但主是此案的审判者，判定不义的一方败诉：因为太阳尚未落山，生理的需要便迫使他去了那地方，他在那里倒下，立即被剥夺了与教会的共融，同时也丧失了生命。真福的\Person{君士坦丁}一听说这事，便惊奇地发现他如此被定为犯伪誓罪。确实，当时所有人都清楚，\Person{尤西比乌}及其同伙的威胁毫无用处，\Person{亚略}的希望也落空了。这也表明，亚略派的狂妄在此处以及天上长子的教会中都遭到了我们救主拒绝其共融。现在，谁不惊讶于看到这些被主定罪的人的不义野心——看到他们维护主已宣布逐出教会（因为祂不容许其创始者进入教会）的异端，并且不惧怕经上所载，反而妄图不可能之事？\BibleQuote{因为万军的上主已决定，谁能取消呢？}\BibleRef{依 14:27} 主所定罪的人，谁能称义呢？然而，让他们为维护自己的想法随意写作吧；但弟兄们，你们作为\BibleQuote{背负上主器皿的}\BibleRef{依 52:11}，并维护教会的教义者，我恳求你们，仔细查考这事；如果他们写出的言论与上面所记\Person{亚略}的言论不同，就谴责他们为伪君子，他们隐藏自己观点的毒素，像蛇一样用口唇的花言巧语谄媚。因为，尽管他们这样写，他们却与那些先前同\Person{亚略}一起被弃绝的人——如\Person{彭塔波利斯的塞昆杜斯}，以及在\Place{亚历山大里亚}被定罪的神职人员——联合，并且写信给在\Place{亚历山大里亚}的他们。但最令人惊讶的是，他们竟然使我们和我们的朋友遭受迫害，尽管最虔诚的皇帝\Person{君士坦丁}曾将我们平安遣返我们的国家和教会，并对他子民的和谐表示关切。但现在他们却使各教会落入这些人之手，从而向所有人证明，正是为了他们的缘故，从起初就策划了针对我们和其他人的整个阴谋。\par

20. \Emph{他们既是\Person{亚略}的朋友，他们温和的话语便是徒劳的。}\par

既然他们的行为如此，他们怎能指望人们相信他们所写的呢？他们若所写下的观点是正统的，早就应从其书目中删除\Person{亚略}的《\WorkTitle{撒里亚}》，并弃绝异端的后裔，即那些\Person{亚略}的门徒，以及与他共同行不虔并受罚的伙伴。但他们并不弃绝这些，所以尽管他们写上千万次，众人也清楚他们的观点并非正统。因此，我们应该警惕，免得在他们的词句外衣下隐藏着欺骗，并勾引某些人离开真信德。如果他们看到自己一帆风顺，竟敢提出\Person{亚略}的观点，我们别无他法，唯有鼓起极大的勇气放胆直言，记起宗徒的预言，这些预言是他为警告我们提防此类异端而写的，我们也应当重述。因为我们知道，正如经上所载：\BibleQuote{在末后的时期，必有人离弃健全的信德，听从欺人的神和魔鬼的训诫，这些训诫偏离真理；} 又如：\BibleQuote{凡愿意在基督里度虔敬生活的，都要受迫害。但恶人和骗子却要变本加厉，迷惑人并受人迷惑。} 但这些事一样也不能胜过我们，也不能\BibleQuote{使我们与基督的爱隔绝}\BibleRef{罗 8:35}，即使异端者以死亡威胁我们。因为我们是基督徒，不是亚略派；但愿那些写这些事的人，没有接受\Person{亚略}的道理！是的，弟兄们，现在极需这样的放胆直言；因为我们所领受的，不是\BibleQuote{再陷于恐惧的奴役之神}\BibleRef{罗 8:15}，而是天主召叫我们\BibleQuote{获得自由}\BibleRef{迦 5:13}。我们若因\Person{亚略}或那些接受并宣扬他观点的人，而毁弃我们从救主透过祂的宗徒所领受的信德，那实在是我们莫大的耻辱，极度的耻辱。在这一带，已有许多人看穿了这些写作者的诡计，甘愿流血以对抗他们的阴谋，尤其当他们听说了你们的坚毅。既然对异端的驳斥已从你们那里发出，而它像洞中之蛇一般被拖出隐藏之处，那么，\Person{黑落德}试图毁灭的那婴孩便保存在你们中间，真理活在你们内，信德在你们中间蓬勃兴旺。\par

21. \Emph{为信德站立等同于殉道。}\par

因此我劝勉你们，将这由尼西亚的教父们所制定的宣认持守在手，满怀热忱，在主内信赖地捍卫它，为各处的弟兄们作榜样，并向他们表明，一场支持真理、反对异端的斗争如今正摆在我们面前，而敌人的诡计是多种多样的。因为殉道者的明证不仅在于拒绝向偶像献香；拒绝否认信德也是良心无愧的光辉见证。不仅那些转向偶像的人被视为异己者而遭弃绝，那些背叛真理的人也是如此。因此，\Person{犹达斯}被剥夺了宗徒职务，不是因为他向偶像献祭，而是因为他显为叛徒；\Person{依默纳约}和\Person{亚历山大}跌倒，不是因为他们转而事奉偶像，而是因为他们\BibleQuote{在信德上倾覆了}（见：\BibleRef{弟前 1:19}）。另一方面，圣祖\Person{亚巴郎}获得了冠冕，不是因为他遭受了死亡，而是因为他忠于天主；\Person{保禄}所提及的其他圣人——\Person{基德红}、\Person{巴辣克}、\Person{三松}、\Person{依弗大}、\Person{达味}和\Person{撒慕尔}，以及其余的人——都不是因流血而成全，而是因信德而成义；直到今日，他们仍是我们钦佩的对象，因为他们为了对主的虔敬甚至随时准备舍生。若可加上一个我国自己的例子，你们知道真福\Person{亚历山大}如何与这异端斗争至死，他虽年迈，却承受了极大的磨难和辛劳，直到高龄时也被接往他祖先那里。此外，还有多少人为了教授道理对抗这不虔而忍受了巨大的劳苦，如今在基督内享受他们宣认的荣耀赏报！因此，我们也要想到这场斗争关乎我们的一切，选择已摆在我们面前：要么否认信德，要么持守信德；让我们也尽心竭力，以在尼西亚所制定的宣认为训导，持守我们所领受的，拒绝新奇，教导我们的子民不要听从\BibleQuote{迷人之神}（见：\BibleRef{弟前 4:1}），而要全心远离那些疯狂的\OriginalTerm{亚略派}{Arians}的不虔，以及\OriginalTerm{梅肋提派}{Meletians}与他们结成的联盟。\par

22. \Emph{卑劣的\OriginalTerm{梅肋提派}{Meletians}与疯狂的\OriginalTerm{亚略派}{Arians}之勾结。}\par

因为你们看得出，他们先前虽彼此不和，如今却像\Person{黑落德}和\Person{般雀}一样，联合起来亵渎我们的主耶稣基督。为此，他们确实该受众人憎恨，因为他们原本因私人原因互相为敌，现在却成了朋友，携手合作，反抗真理，不虔事奉天主。不仅如此，他们为达成各自的目的，不惜做出或忍受任何与他们原则相悖的事；\OriginalTerm{梅肋提派}{Meletians}为了超越他人和疯狂地贪财，疯狂的\OriginalTerm{亚略派}{Arians}则为着自己的不虔。这样，通过这种勾结，他们得以互相帮助完成恶毒的阴谋，梅肋提派披上了亚略派的不虔，亚略派也以自己的邪恶助长了他们的卑鄙，从而把各自的罪行混合在一起，如同\Place{巴比伦}的杯（见：\BibleRef{默 18:6}），为要谋害我们主耶稣基督的正统敬拜者。梅肋提派的邪恶与虚假，其实此前已为人人皆知；亚略派的不虔与无神的异端也早为各地众人所知；因为他们存在的时间并不短。前者于五十五年前成为裂教者，而后者被宣判为异端者已有三十六年，且经整个\OriginalTerm{大公会议}{Ecumenic Council}的审判被摈弃于教会之外。但通过他们当前的所作所为，最终甚至对那些看似公开支持他们的人证明了，他们对我及其他正统主教施展阴谋，从一开始就单单是为了推进自己那份不虔的异端。看哪，\Person{欧瑟比}及其同伙长久以来所谋取的大业，如今已经实现。他们导致了各教会从我们手中被夺去，任意放逐不愿与他们共融的主教和长老，并将那脱离他们的人逐出教会，又将教会交于那早已被定罪的\OriginalTerm{亚略派}{Arians}之手，如此，借着梅肋提派的伪善，他们可以肆无忌惮地在其中散布不虔的言论，并如他们所想，为那在他们中间撒下这异端种子的\OriginalTerm{反基督}{Antichrist}预备那欺谎之路。\par

23. \Emph{结论。}\par

然而，任他们如此梦想和想像虚空之事吧。我们知道，当我们仁慈的皇帝听闻此事时，他必会制止他们的邪恶，他们不会长久持续，但正如圣经所言：\BibleQuote{不虔之人的心必快快衰残。}但我们如经上所载，要\BibleQuote{穿上圣经的话，}并抵抗他们这些背教者，他们要在主的殿中建立狂热崇拜。我们不要怕身体的死亡，也不要效法他们的行径；但要让真理之言在一切事上优先。正如你们都知道的，我们先前也被\Person{优西比乌}和他的同伙要求，要么接受他们的不虔敬，要么等着他们敌对；但我们不愿与他们同流合污，宁可受他们迫害，也不效法\Person{犹达斯}的行径。他们确实行了他们威胁的事；因为效法\Person{依则贝耳}的样式，他们利用背信弃义的\Person{梅勒提乌}派来帮助他们，知道后者如何抗拒了有福的殉道者\Person{伯多禄}，在他之后有大\Person{阿基拉斯}，然后是令人怀念的\Person{亚历山大}，好使\Person{梅勒提乌}派在这类事上练达，也照样诬告我们，任凭\Person{优西比乌}和其同伙提供迫害和谋害我的机会。这就是他们所渴求的；他们直到今日还想流我的血。但我对这些事毫不挂心；因为我知道并深信，忍耐的人必从我们的救主那里获得赏报；你们也是如此，如果你们像教父们那样忍耐，在人民面前立下榜样，推翻这些不虔之人奇怪而诡诈的计谋，你们就能夸耀地说：\BibleQuote{我保持了信德}\BibleRef{弟后4:7}，并且你们必领受\BibleQuote{生命的冠冕}，就是天主\BibleQuote{向爱祂的人所应许的}\BibleRef{雅1:12}。愿天主准许我也同你们一起承受那应许，这应许不仅赐给了\Person{保禄}，也赐给了所有\BibleQuote{爱慕祂显现的人}\BibleRef{弟后4:8}——我们的主、救主、天主、普世的君王耶稣基督；愿光荣与权能，借着祂，在圣神内，归于圣父，从今直到永远，世世无穷。阿们。\par

\SourceRecord{Translated by M. Atkinson and Archibald Robertson.；Nicene and Post-Nicene Fathers, Second Series；Vol. 4.；Edited by Philip Schaff and Henry Wace.；Buffalo, NY: Christian Literature Publishing Co.,；1892.；<http://www.newadvent.org/fathers/2812.htm>.}
